\documentclass{article}
\usepackage[utf8]{inputenc}
\usepackage[spanish]{babel}
\usepackage{graphicx}
\usepackage{float}
\usepackage{adjustbox}
\usepackage{subcaption}    
\usepackage{geometry}
\usepackage{tabulary}

\geometry{
   a4paper,
   total={170mm,257mm},
   left=35mm,
   right=35mm,
   top=30mm,
   bottom=30mm,
}
 
\title{Propuesta de acuerdo del proyecto}
\author{Ingeniería web\\\\Grupo 13}
\date{Octubre 2021}
\setlength{\parindent}{2em}
\setlength{\parskip}{1em}
 
\begin{document}
 
\maketitle
 
\begin{figure}[H]
    \captionsetup[subfigure]{justification=centering}
     \begin{subfigure}[b]{0.49\textwidth}
        \centering
         \includegraphics[width=0.5\textwidth]{../images/Andoni.jpeg}
         \caption*{\textbf{Lider}: Salcedo Navarro, Andoni (785649)}
     \end{subfigure}
     \hfill
     \begin{subfigure}[b]{0.49\textwidth}
        \centering
        \includegraphics[width=0.5\textwidth]{../images/Tomas.jpeg}
        \caption*{Pelayo Benedet, Tomás (779691)}
     \end{subfigure}
     \hfill
     \begin{subfigure}[b]{0.49\textwidth}
        \centering
        \includegraphics[width=0.5\textwidth]{../images/Andoni.jpeg}
        \caption*{Subías Rodríguez, Rubén (759406)}
     \end{subfigure}
     \hfill
     \begin{subfigure}[b]{0.49\textwidth}
        \centering
        \includegraphics[width=0.5\textwidth]{../images/Andoni.jpeg}
        \caption*{Velasco Calvo, Isaac (758986)}
     \end{subfigure}

    \end{figure}

\pagebreak
 
\section*{Objetivos}

Se han analizado las caracteristicas del equipo de trabajo y los requerimientos que son necesarios para afrontar el proyecto y se ha llegado a la conclusión de que los objetivos del equipo de cara al proyecto son los siguientes:

\begin{itemize}
    \item O1. Ser capaz de identificar las necesidades para desarrollar una aplicación Web que se pueda despliegar en producción.
    \item O2. Ser capaz de adaptar las nuevas funcionalidades propuestas y desplegarlas en el sistemas web ya funcional.
    \item O3. Comprender como funcionan las pasarelas publicitarias para poder monetizar la aplicación.
    \item O4. Comprensión de las arquitectura detras de los  canales de comunicación multicast.
    \item O5. Entender de como gestinar el almacenamiento y encriptado de datos sensibles.
    \item O6. Aprender a utilizar los conceptos de UX/UI para que la usabilidad de la aplicación sea la adecuada para el usuario final.
\end{itemize}

\section*{Funcionalidades}

\begin{itemize}
    \item F1. La aplicación implementará un chat grupal en el que se podrán compartir enlaces con otras personas que usen la aplicación (O1,O2,O4). Con el objetivo de aprender como implementar canales de comunicación bidireccionales.
    \item F2. La aplicación soportará un inicio de sesión (O1,O2,O5,O6). Con el objetivo de que cada usuario pueda acceder a sus uri's creadas para poder modificarlas, compartirlas, añadir, eliminar, monetizar, con la finalidad de entender como se realiza el control de sesión y gestión de contraseñas en un nuevo framework.
    \item F3. Se incluye un sistema de redirección para que el usuario pueda monetizar el número de clicks sobre sus uri's (O1,O3). Con el objetivo de entender como se crea un modelo de negocio real a partir de una aplicación.
    \item F4. La aplicación verificará que la uri que el usuario quiere crear es alcanzable (O2,O5). El objetivo es aprender como se puede comprobar si una redirección es valida, para evitar redirecciones dañinas o inexistentes.
\end{itemize}

\begin{table}[hbtp]
    \footnotesize
    \centering
    \settowidth\tymin{\textbf{Activities}}
    \setlength\extrarowheight{5pt}
    \begin{tabulary}{\textwidth}{ C C L C}
        \textbf{ID} & 
        \textbf{Peso} & 
        \textbf{Funcionalidad} & 
        \textbf{Objetivos}\\
    \hline
    \hline
        F1 & 
        15 & 
        La aplicación implementará un chat grupal en el que se podrán compartir enlaces con otras personas que usen la aplicación. &
        O1, O2, O4\\
    \hline  
        F2 &
        10 &
        La aplicación soportará un inicio de sesión. &
        O1, O2, O5, O6\\
    \hline
        F3 &
        25 &
        Se incluye un sistema de redirección para que el usuario pueda monetizar el número de clicks sobre sus uri's. &
        O1, O3\\
    \hline
        F4 &
        10 &
        La aplicación verificará que la uri que el usuario que quiere crear es alcanzable. &
        O2, O5

    \end{tabulary}
    \caption{Tabla de reparto de pesos.}
\end{table}

\pagebreak

\section*{Orientaciones sobre la descripción de funcionalidades}


\begin{table}[hbtp]
    \footnotesize
    \centering
    \settowidth\tymin{\textbf{Funcionalidad}}
    \setlength\extrarowheight{5pt}
    \begin{tabulary}{\textwidth}{ C L }
        \textbf{F1} & La aplicación implementará un chat grupal en el que se podrán compartir enlaces con otras personas que usen la aplicación.
        \\
    \hline
    
    Cumplimiento & Los usuarios pueden entrar a un chat, en el pueden interactuar entre ellos, todos los usuarios participantes en el chat reciben todos los mensajes, en el chat se pueden compartir links de los enlaces recortados de cada usuario.\\
        
    Escalabilidad & Los mensajes deben llegar en el orden de envío se asegura de que  \\

    Profesionalidad & dfasdfadsasdf \\

    \end{tabulary}
\end{table}

\begin{table}[hbtp]
    \footnotesize
    \centering
    \settowidth\tymin{\textbf{Funcionalidad}}
    \setlength\extrarowheight{5pt}
    \begin{tabulary}{\textwidth}{ C L }
        \textbf{F2} & La aplicación soportará un inicio de sesión.
        \\
    \hline
    
    Funcionalidad & adfadfasdfasdf\\
        
    Escalabilidad & fafadfa \\

    Profesionalidad & dfasdfadsasdf \\

    \end{tabulary}
\end{table}

\begin{table}[hbtp]
    \footnotesize
    \centering
    \settowidth\tymin{\textbf{Funcionalidad}}
    \setlength\extrarowheight{5pt}
    \begin{tabulary}{\textwidth}{ C L }
        \textbf{F3} & Se incluye un sistema de redirección para que el usuario pueda monetizar el número de clicks sobre sus uri's.
        \\
    \hline
    
    Funcionalidad & adfadfasdfasdf\\
        
    Escalabilidad & fafadfa \\

    Profesionalidad & dfasdfadsasdf \\

    \end{tabulary}
\end{table}

\begin{table}[hbtp]
    \footnotesize
    \centering
    \settowidth\tymin{\textbf{Funcionalidad}}
    \setlength\extrarowheight{5pt}
    \begin{tabulary}{\textwidth}{ C L }
        \textbf{F4} & La aplicación verificará que la uri que el usuario quiere crear es alcanzable. 
        \\
    \hline
    
    Funcionalidad & Solo se puede crear una URI acortada si se verifica que una petición HTTP GET
    a la URI original devuelve una respuesta con estado 200.\\
        
    Escalabilidad & crea la URI acortada pero no se activa; en un hilo de ejecución separado se
    verifica que la URI original es alcanzable mediante una petición http GET que debe devolver
    una respuesta con estado 200 en un tiempo máximo configurable; si es alcanzable la URI
    acortada se activa, si no lo es la URI acortada se marca como inválida. \\

    Profesionalidad & toda la funcionalidad está cubierta por pruebas automáticas de unidad u de
    integración que comprueban al menos dos casos (URI alcanzable, URI no alcanzable), el API
    relacionado está documentado utilizando Open API, hay una consola Web interactiva que
    visualiza dicha documentación y que permite probar dicha API y hay documentación en
    AsciiDoc explicando el funcionamiento del código que ha sido convertida a HTML5 con
    Asciidoctor.
    \end{tabulary}
    
\end{table}
 
\end{document}
